% drafts/07-extensions.tex — Extension to Bosonic and Mixed Systems
%
% Target length: ~2 pages (two-column PRA)
% Subsections: 7.1 Design principle, 7.2 Bosonic algebra,
%              7.3 Bosonic-to-qubit encodings, 7.4 Mixed systems

The framework developed in \cref{sec:index-set,sec:tree-encoding} is
specific to fermions: the encoding maps are designed to satisfy the
CAR.  We now show that the \emph{same algebraic pipeline} — product
types, normal ordering, coefficient tracking — extends to bosonic
systems (canonical commutation relations, CCR) and to mixed
fermion--boson Hamiltonians, requiring only a change of the combining
algebra.

%%====================================================================
\subsection{The combining algebra abstraction}
\label{sec:ext-algebra}

The normal-ordering algorithm for ladder operators works by moving
creation operators to the left of annihilation operators, applying
the fundamental relation each time a swap is needed.  The only
difference between particle statistics is the relation itself:
\begin{align}
  \text{Fermionic (CAR):}\quad
  \acomm{\an{i}}{\ad{j}} &= \delta_{ij},
  \label{eq:car-combine} \\
  \text{Bosonic (CCR):}\quad
  \comm{b_i}{b_j^\dagger} &= \delta_{ij}.
  \label{eq:ccr-combine}
\end{align}

We abstract this into a single interface parameterized by the
commutation algebra:
\begin{definition}[Combining algebra]
\label{def:combining}
A \emph{combining algebra} $\mathcal{A}$ is a function
\begin{equation}
  \textsc{Combine}: P \times C \to P[],
\end{equation}
where $P$ is a product of ladder operators and $C$ is the next
operator to append.  When a $b_j$ must pass a $b_k^\dagger$:
\begin{itemize}
  \item If $j \neq k$: return one term, with sign $-1$ (fermionic)
    or $+1$ (bosonic).
  \item If $j = k$: return \emph{two} terms — the commutator/
    anticommutator identity contribution (a shorter product) and the
    reordered product with sign $\mp 1$.
\end{itemize}
\end{definition}

The normal-ordering algorithm is identical for both statistics; only
the \textsc{Combine} function differs.  This is a concrete instance
of the ``algebra as parameter'' pattern: the pipeline is a functor
from combining algebras to normal-ordered expressions.

%%====================================================================
\subsection{Bosonic normal ordering}
\label{sec:ext-bosonic}

For bosonic modes, swapping $b_j$ past $b_k^\dagger$ with $j = k$
produces two terms:
\begin{equation}
  b_j b_j^\dagger = b_j^\dagger b_j + \mathbb{1},
  \label{eq:ccr-swap}
\end{equation}
with \emph{no} sign change on the reordered term (contrast the
fermionic case, where $a_j a_j^\dagger = -a_j^\dagger a_j +
\mathbb{1}$ effectively, via the anticommutator).  For $j \neq k$,
bosonic operators commute: $b_j b_k^\dagger = b_k^\dagger b_j$,
again with no sign change.

The normal-ordering of a generic bosonic product proceeds by the same
bubble-sort algorithm as the fermionic case, substituting the CCR
\textsc{Combine} function.  The result is a sum of normal-ordered
terms $\sum_m c_m \prod_\ell (b^\dagger)^{n_\ell} b^{m_\ell}$
with exact coefficients.

%%====================================================================
\subsection{Bosonic-to-qubit encodings}
\label{sec:ext-boson-qubit}

Unlike fermions (where each mode maps to one qubit), bosonic modes
have unbounded occupation and must be \emph{truncated} to a maximum
occupation number $d$ before mapping to qubits.  Given $d$, each
bosonic mode occupies $q(d)$ qubits and the operators are mapped:

\begin{description}
  \item[Unary encoding ($q = d$ qubits):] Each Fock state $\ket{n}$
    is represented by a one-hot basis state.  The creation operator
    is decomposed as transitions between adjacent occupation levels:
    \begin{equation}
      b^\dagger = \sum_{n=0}^{d-2} \sqrt{n+1}\; \sigma^+_{n+1}\,\sigma^-_n,
      \label{eq:unary}
    \end{equation}
    where $\sigma^\pm = \frac{1}{2}(X \mp iY)$.
    Each transition has Pauli weight~$\le 2$, giving
    $4(d-1)$ Pauli terms per mode.

  \item[Binary encoding ($q = \lceil\log_2 d\rceil$ qubits):]
    The truncated bosonic operator is represented as a $d \times d$
    matrix $(b^\dagger)_{n+1,n} = \sqrt{n+1}$, embedded into
    $2^q \times 2^q$, and decomposed via the Pauli trace formula:
    \begin{equation}
      b^\dagger = \sum_{P \in \{I,X,Y,Z\}^{\otimes q}}
        \frac{\mathrm{Tr}(P \cdot b^\dagger)}{2^q}\; P.
      \label{eq:binary}
    \end{equation}
    Maximum Pauli weight is $q$.

  \item[Gray-code encoding ($q = \lceil\log_2 d\rceil$ qubits):]
    Similar to binary, but with the basis ordering permuted by
    the Gray code $n \mapsto n \oplus (n \gg 1)$.  Consecutive
    Fock states differ in exactly one qubit, improving the average
    Pauli weight at the cost of identical maximum weight.
\end{description}

\begin{remark}
The bosonic encodings introduce irrational coefficients
($\sqrt{n+1}$) and use the matrix-decomposition trace formula, unlike
the fermionic path which is entirely symbolic.  Phase exactness
(\cref{sec:exactness}) does not extend to the bosonic case; this is an
inherent feature of continuous-variable truncation, not a limitation
of the pipeline.
\end{remark}

%%====================================================================
\subsection{Mixed fermion--boson systems}
\label{sec:ext-mixed}

Electron--phonon and polariton--cavity Hamiltonians contain both
fermionic and bosonic operators.  The standard approach is to treat
cross-sector operators as commuting: a fermionic $\ad{j}$ commutes
with a bosonic $b_k^\dagger$ because they act on distinct Hilbert
spaces.

Our pipeline handles mixed systems by:
\begin{enumerate}
  \item \textbf{Sector tagging.}  Each operator carries a label
    $\mathrm{Fermionic}$ or $\mathrm{Bosonic}$, so
    $\ad{j}^{(F)}$ and $b_k^{\dagger\,(B)}$ are syntactically
    distinct.
  \item \textbf{Block reordering.}  Within each product, all
    fermionic operators are moved to the left and bosonic operators
    to the right.  Cross-sector swaps introduce no sign change
    (they commute).
  \item \textbf{Sector-wise normal ordering.}  The fermionic
    sub-sequence is normal-ordered using the CAR algebra; the
    bosonic sub-sequence is normal-ordered using the CCR algebra.
  \item \textbf{Cartesian recombination.}  The fermionic and bosonic
    normal-ordered expansions are combined by distributing
    coefficients multiplicatively.
\end{enumerate}

The result is a sum of fully normal-ordered mixed-sector products
with coefficient tracking inherited from both algebras.  This
four-step algorithm is implemented generically: it requires only
the \textsc{Combine} function from each sector's algebra, plus the
commutativity assumption for cross-sector swaps.

\begin{remark}
The qubit encoding of a mixed system is obtained by composing the
fermionic encoding (\cref{sec:index-set} or \cref{sec:tree-encoding})
with a bosonic encoding (\cref{sec:ext-boson-qubit}) on disjoint
qubit registers.  The two registers are tensored together in the
final Pauli representation.
\end{remark}
